\subsection{外的導分の跳躍を有限個の方程式へ縮約する}

量子VKの三十五変数hull $R$ 上では、相対Hochschild複体を単元消去すると
次数一から二への一列 $a:R\to R^{35}$ だけが非零の微分として残る。
次数零は $R$、次数三は $R^{39}$ である。
結合積スキームが六百五十一次元で滑らかなため、$d^2$ の階数は九百三十三で一定となる。
この点が、特殊ファイバーの階数だけから残余微分の消滅を推論しないために必要である。

$I=(a_i)=\operatorname{Fitt}_{34}H^2(C_R^\bullet)$ は、外的導分の持上げが存在する
閉形式部分 $Z$ を定める。列の一次項は $-L_\delta$、$\delta=\log(1+N)$ である。
従って $T_0Z=\ker L_\delta$、次元は $j=0$ で四、$j=1$ で六。
$Z$ の外では $HH=(1,0,34,39)$ となり、通常のhullは三十四変数で滑らかである。
$R$ 上の相対 $HH^1$ 自体は零なので、特殊化との非可換性を見落とさない。
跳躍は複体の導来基底変換で保持される。

独立な一次方程式を消すと、$Z$ の全環は
\[
 \mathbf C[[v_1,\ldots,v_k]]/(f_1,\ldots,f_k),\qquad f_i\in(v)^2,
 \qquad k=4\text{ または }6
\]
に縮約される。二次項と、後述の全次数の族を組み合わせて高次構造を限定できる。
$j=0$ では全て零、$j=1$ では一つの平方 $v_6^2$ が生成する。
この平方が消える超平面は、内在的には $\ker L_\delta\cap\operatorname{im}L_\delta$ である。
従って $j=1$ の外的導分跳躍部分は原点で特異で、次元は正確に五であり、滑らかな五次元の閉形式部分を持つ。
二次項が平方であることだけから全環の非被約性を推論しない。

再利用する二次障害の公式は、$u\in\ker L_\delta$ に対して
\[
 d^1\eta=L_\delta u,\qquad d^2v=u\circ u,
 \qquad
 \mathfrak q(u)=[L_\delta v+L_\eta u]\bmod\operatorname{im}L_\delta
\]
である。代数変形の無障害性と、導分を同時に持ち上げる条件は異なる。
独立な検算は $2556\times2204$ のblock行列
$\bigl(\begin{smallmatrix}d^2&0\\L_\delta&-d^1\end{smallmatrix}\bigr)$ で
二つの持上げ条件を同時に解く方法であり、第二次の積補正の選択に依存しない。

ここで保持するのは非零外的導分の存在であり、指定した $\mathbf G_a$ 作用や
局所冪零性を含む同変変形ではない。
全対称性保持部分は $j=0$ で滑らかな四次元、$j=1$ では滑らかな五次元部分の二重構造となる。
標準補完線と五次元接平面の位置関係も決定できる。
純粋三変数部分への還元と全二重構造の証明は後述する。
